\subsection{不变因子的计算}

命题4. --- 设 A 为主理想整环，L 为具有有限基 \((u_{j})\) \((1 \leqslant j \leqslant k)\) 的自由 A-模，设 M 是 L 的一个子模，设 \((x_{i})\) 是 M 的一个生成元组，并设 \(A\alpha_{i}\) \((1 \leqslant i \leqslant n)\) 是 M 关于 L 的不变因子。则对 \(1 \leqslant m \leqslant n\)，乘积 \(\delta_{m} = \alpha_{1} \ldots \alpha_{m}\) 是以诸 \(x_{i}\) 关于基 \((u_{j})\) 的坐标向量为列的矩阵的 m 阶子式的一个最大公因子。

由定理 1 显然可知 \(M \subset \alpha_{1}L\)，因而 M 中任何元素的坐标都是 \(\alpha_{1}\) 的倍数。另一方面，存在 M 的一个元素 x，使 \(\alpha_{1}\) 是 x 在 L 中的一个内容。把 x 表示成诸 \(x_{t}\) 的线性组合，便知 \(\alpha_{1}\) 属于由诸 \(x_{t}\) 的坐标生成的理想。由于这些坐标都是 \(\alpha_{1}\) 的倍数，于是 \(\alpha_{1}\) 确是它们的最大公因子，从而我们的断言在 m = 1 的情形得证。

对于一般的 \(m\)，考虑 M 的第 \(m\) 个外幂 \(\Lambda\) M（III，第 507 页）。在定理 1 的记号下，存在 M 的一组基 \((a_i)\)，其中 \(a_i = \alpha_i e_i\)（\(1 \leqslant i \leqslant n\)）；于是 \(\Lambda\) M 有一个由诸元素 \(a_{i_1} \wedge \ldots \wedge a_{i_m}\) 组成的基，其中 \((i_1, \ldots, i_m)\) 遍历由 \(m\) 个取自 \((1, n)\) 的元素组成的严格递增序列的集合。而诸元素 \(e_{i_1} \wedge \ldots \wedge e_{i_m}\) 属于一组基 \(B_m\)（\(\bigwedge^m L\) 的基）。因此从 \(\bigwedge^m M\) 到 \(\bigwedge^m L\) 的典范映射是从 \(\bigwedge^m M\) 到 \(\bigwedge^m L\) 的某个子模上的一个同构，该子模以诸元素 \((\alpha_{i_1} \ldots \alpha_{i_m}) e_{i_1} \wedge \ldots \wedge e_{i_m}\) 为基，我们把这个子模与 \(\bigwedge^m M\) 等同起来。由于 \(\alpha_j\) 是 \(\alpha_k\) 的倍数（对 \(j \geqslant k\)），诸元素 \(\alpha_{i_1} \ldots \alpha_{i_m}\) 都是 \(\delta_m = \alpha_1 \ldots \alpha_m\) 的倍数，且其中之一等于 \(\delta_m\)；于是 \(\delta_m\) 是关于基 \(B_m\) 取的 \(\bigwedge^m M\) 的某个生成元组中诸元素的坐标的一个最大公因子。论证的第一部分表明，此时 \(\delta_m\) 是 \(\bigwedge^m M\) 的任意生成元组关于 \(\bigwedge^m L\) 的任意基的坐标集的一个最大公因子。取 \(\bigwedge^m L\) 的由 \((u_j)\)（即 \(L\) 的基）诱导的基，以及 \(\bigwedge^m M\) 的由诸 \((x_i)\) 的外积组成的生成元组，则这些外积的坐标借行列式的表达式（III，第 528 页，命题 9）便给出所述结果。

